\begin{longtblr}[
  label={tab:antecedentes_intro},
  caption={Antecedentes que delimitan el alcance de la revisión sistemática.}
]{
  colspec={X[0.8,c,m] X[1.8,l,m] X[2.4,l,m] X[2.5,l,m]},
  width =\fulllength,
  rowhead=1,
  row{1}={font=\bfseries}
}
\toprule
Referencia & Tipo/aporte principal & Relación con este estudio & Limitación frente a esta SLR \\
\midrule
Priti y Kumar (2022) \cite{ART01} & Revisión crítica de modelos AQI (1960--2021). & Delimita el problema de comparabilidad entre índices y la heterogeneidad metodológica del campo. & No aborda inferencia difusa ni monitoreo en tiempo real como foco de análisis. \\
Suman (2020) \cite{ART02} & Revisión de métodos para interpretar el estado de la calidad del aire mediante AQI. & Refuerza la función comunicativa del AQI y su relación con la interpretación del riesgo. & No analiza decisiones de diseño difuso ni arquitecturas de monitoreo inteligente. \\
Murena (2004) \cite{PROB08} & Desarrollo y aplicación de un índice de contaminación en contexto urbano. & Acerca la discusión del AQI a un caso de implementación sobre área urbana amplia. & No incorpora inferencia difusa ni una lectura comparativa de plataformas en tiempo real. \\
Seibert et al. (2024) \cite{ART04} & Clasificación difusa de calidad del aire. & Representa la línea de uso de reglas difusas para clasificar estados de calidad del aire. & Se concentra en una tarea específica de clasificación y no en una síntesis de configuraciones de diseño. \\
Andrianto et al. (2024) \cite{ART06} & Monitoreo IoT con lógica difusa Mamdani e interpretación tipo ISPU/AQI. & Aporta un antecedente cercano a la integración entre sensores, índice e interpretación en tiempo real. & Corresponde a una implementación puntual y no permite por sí sola comparar variantes metodológicas del campo. \\
Lingaraj et al. (2025) \cite{ART08} & Arquitectura WSN optimizada para monitoreo de aire en smart cities. & Introduce la dimensión de operación distribuida y eficiencia de red en despliegues urbanos. & Privilegia la arquitectura de red y no articula de forma transversal variables, diseño difuso y evaluación del sistema. \\
\bottomrule
\end{longtblr}
